\begin{trapbox}
\textbf{易错点 1：把“相切”忘成“任意直线”}\par
$d_1d_2=b^2$ 只对\textbf{椭圆的切线}成立。比如直线 $y=0$ 经过两个焦点，到两焦点的距离都为 $0$，乘积不是 $b^2$；该直线与椭圆相交于两个长轴端点，\textbf{不是切线}。看到一般式 $Ax+By=D$，先检查 $D^2=a^2A^2+b^2B^2$。

\medskip
\textbf{易错点 2：用 $y=kx+m$ 不加说明地覆盖全部切线}\par
该写法漏掉 $x=\pm a$。竖直切线对应“斜率不存在”，不是 $k=\infty$ 可以直接代入普通公式。建议优先用 $Ax+By=D$ 或单位法向量式；若用斜率法，请补上竖直切线。

\medskip
\textbf{易错点 3：把切点当成焦点的垂足}\par
切点 $P$ 满足椭圆方程 $x^2/a^2+y^2/b^2=1$；垂足 $H_i$ 满足\textbf{主圆方程} $x^2+y^2=a^2$。除某些特殊位置外，两者不是同一点。尤其切线 $y=b$ 的右焦点垂足是 $(c,b)$，不是 $(0,b)$。

\medskip
\textbf{易错点 4：只证明“垂足在圆上”，就声称“完整轨迹是整个圆”}\par
证明 $OH=a$ 仅说明垂足集合\textbf{包含于}主圆；若题目问完整轨迹，必须再证明主圆上任意点都可作某条切线的垂足。逆向构造：取 $H=(p,q)$，$p^2+q^2=a^2$，固定 $F_2(c,0)$，作 $l:(p-c)x+qy=a^2-cp$，验算切线判定式即可。

\medskip
\textbf{易错点 5：混淆主圆、辅圆与焦点圆}\par
这里的主圆（长轴辅助圆）是 $x^2+y^2=a^2$，圆心为椭圆中心 $O$；不是半径 $b$ 的圆，也不是以焦点为圆心、半径 $2a$ 的反射用圆。几何反射中出现的“以另一焦点为圆心、半径 $2a$ 的圆”只用于帮助理解，不等于垂足轨迹。

\medskip
\textbf{易错点 6：距离必须非负，切线两焦点位于同侧}\par
点到直线距离公式必须加绝对值：$d(F,l)=|Ax_F+By_F-D|/\sqrt{A^2+B^2}$。只有将切线写成\textbf{外法向单位式} $ux+vy=h$（$h>0$）并先证明 $h>|cu|$，才可直接写 $d_1=h+cu$、$d_2=h-cu$ 而不加绝对值。

\medskip
\textbf{易错点 7：将距离积错记为 $a^2$ 或 $c^2$}\par
正确配对：\textbf{焦点距离积为 $b^2$；垂足主圆半径为 $a$；焦距的一半为 $c=\sqrt{a^2-b^2}$。} 可以用竖直切线 $x=a$ 快速检验：$(a-c)(a+c)=a^2-c^2=b^2$。

\medskip
\textbf{易错点 8：两焦点距与中心距关系的适用前提}\par
$d_1+d_2=2d_0$ 是\textbf{切线}的性质，因为中心与两个焦点同在切线的内侧；对横穿焦点连线的普通直线，不能直接使用此和式。

\medskip
\textbf{易错点 9：椭圆方向改变后死记坐标形式}\par
本章规定长轴在 $x$ 轴上，$a>b>0$。若长轴改在 $y$ 轴，焦点应写作 $(0,\pm c)$，垂足依然在以中心为圆心、半径为\textbf{长半轴 $a$} 的圆上，焦点距离积仍为\textbf{短半轴的平方 $b^2$}。不要把焦点坐标照抄。

\medskip
\textbf{易错点 10：信息不足却强行求唯一方程}\par
“焦点到切线距离之积为 $4$”只能推出 $b^2=4$。即使再说椭圆经过 $(0,2)$，仍无法确定 $a$；需要另给离心率、焦距或长轴长度之类的独立条件。
\end{trapbox}
